% !TITLE 临路歌
% !AUTHOR 李白
% !DYNASTY 唐
% !ORDER 17

\begin{poem}
大鹏飞兮振八裔\textsuperscript{1}，中天摧兮力不济\textsuperscript{2}。
余风激兮万世\textsuperscript{3}，游扶桑兮挂石袂\textsuperscript{4}。
后人得之传此，仲尼亡兮谁为出涕\textsuperscript{5}。
\end{poem}

\begin{annotations}
\notetext{1}{八裔：八个方向的荒远之地，即八方极远处。大鹏展翅振动了八方。}
\notetext{2}{中天摧兮力不济：在半空中翅膀被摧折，力气接不上了。李白以此自况一生的壮志未酬。}
\notetext{3}{余风激兮万世：大鹏余下的风力依然能激荡万世。即使失败了，影响还会留存。}
\notetext{4}{扶桑：神话中太阳升起处的神树。石袂：左袖。大鹏曾在扶桑树上栖息，左袖挂在树枝上。}
\notetext{5}{仲尼亡兮谁为出涕：孔子（字仲尼）已经死了，谁还能为大鹏的坠落流泪？典故：鲁人捕获麒麟，孔子见之而泣。李白以麒麟自比，感慨世上再无知音。}
\end{annotations}

\begin{appreciation}
这是李白的绝笔，写于他生命最后的年头，也是中国文学史上最沉重的遗言诗之一。诗只有六句，说的是一只大鹏。

大鹏飞兮振八裔。八裔，是八方的边远之地。大鹏展翅，振动八方。这个意象李白用了一辈子——还记得《上李邕》吗？“大鹏一日同风起，扶摇直上九万里”，那是二十多岁的李白，觉得自己一扇翅膀就能到九万里高。大鹏出自《庄子·逍遥游》，庄周写它“怒而飞，其翼若垂天之云”，本来是形容自由的极限，被李白拿来当自己的名字。

中天摧兮力不济。飞在半空，翅膀折了，力气接不上了。六个字，把自己一生的结局写完了。当年那个扶摇直上的年轻人，终究没飞到九万里，在人生的中途跌了下来。这不是谦虚——到了临终，没必要谦虚，他只是说了实话。

余风激兮万世。大鹏掉下来，带起的风还在。这句是李白留给世界的信心：我这一生没飞成，但翅膀扇过的风，能吹到后世。游扶桑兮挂石袂。扶桑是神话里太阳升起的地方，大鹏曾在太阳边栖息，衣袖挂在了树枝上。到过最高的地方，即使被挂住，也算数。

最后一句最沉：后人得之传此，仲尼亡兮谁为出涕。传说鲁人捕到一头麒麟，孔子见了落下泪来——麒麟是祥瑞之物，不该在乱世被猎杀，此后孔子绝笔不写《春秋》。李白把麒麟换成大鹏，问：孔子已经不在了，还有谁为我这只坠落的大鹏流泪？一个一生自负“天生我材必有用”的人，到最后在乎的不是官位、不是功名，是“有没有人懂我”。

李白一辈子都在写自己。《南陵别儿童入京》里，他出门去长安，笑得那么大声：“仰天大笑出门去，我辈岂是蓬蒿人”。那时他不知道，长安终究不是他的地方。到了《临路歌》，他不再笑，只把一生收进一只折翼的大鹏。从“扶摇直上九万里”到“中天摧兮力不济”，一个人的一生，两首诗就装完了。

读完这首诗，哥哥想对你说的话其实很简单：李白飞了一辈子也没飞到想去的地方，但他从来没有后悔张开过翅膀。你以后也会有很多想做的事，做成做不成都不打紧，要紧的是肯不肯扑腾。至于飞多高，那是风的事；肯不肯飞，是你自己的事。
\end{appreciation}
